% !TEX program = lualatex
% !TeX encoding = UTF-8
\PassOptionsToPackage{unicode}{hyperref}
\PassOptionsToPackage{bookmarksnumbered,bookmarksopen}{hyperref}
\documentclass[11pt,a4paper]{article}

% ============================================================
% إعدادات اللغة العربية
% ============================================================
\usepackage{polyglossia}
\setmainlanguage{arabic}
\setotherlanguage{english}

% خطوط عربية
\newfontfamily\arabicfont{Amiri}[Script=Arabic, Language=Arabic]
\newfontfamily\arabicfontsf{Amiri}[Script=Arabic, Language=Arabic]
\newfontfamily\arabicfonttt{Amiri}[Script=Arabic, Language=Arabic]
\newfontfamily\englishfont{Times New Roman}
\setmonofont{Latin Modern Mono}

% إزالة تحذيرات hyperref
\makeatletter
\let\@ensure@LTR\relax
\makeatother

% حزم الرياضيات والتنسيق
\usepackage{amsmath,amssymb,amsthm,mathtools,bm}
\usepackage{geometry}
\usepackage{enumitem}
\usepackage{siunitx}
\usepackage{booktabs}
\usepackage{array}
\usepackage{slashed}
\usepackage{graphicx}
\usepackage{caption}
\usepackage{subcaption}
\usepackage{braket}
\usepackage{tensor}
\usepackage{makecell}
\usepackage[most]{tcolorbox}
\usepackage{pgfplots}
\usepackage{float}
\usepackage{tikz}
\usepackage{multicol}
\usepackage{lipsum}
\usepackage{microtype}
\usepackage{hyperref}
\pgfplotsset{compat=1.18}
\usetikzlibrary{arrows.meta,positioning,shapes.geometric,fit,calc,
	decorations.pathreplacing,patterns,quotes,angles,
	shadows,decorations.markings}

\geometry{margin=1in}
\setlength{\emergencystretch}{3em}
\sloppy

% بيئات النظريات
\newtheorem{theorem}{نظرية}
\newtheorem{lemma}{ليما}
\newtheorem{axiom}{بديهية}
\newtheorem{remark}{ملاحظة}
\newtheorem{proposition}{اقتراح}
\newtheorem{definition}{تعريف}
\newtheorem{assumption}{افتراض}
\newtheorem{corollary}{نتيجة طبيعية}

% بيانات PDF
\hypersetup{
	pdftitle={الملحق Ferra1.2: الثوابت التناسقية الشاملة للأطوار المرتبة وغير المرتبة},
	pdfauthor={أليكسي ف. ياكوشيف},
	pdfsubject={YUCT, ثوابت تناسقية شاملة, طور مرتب, طور غير مرتب, مغناطيسية حديدية, تحول طوري},
	pdfkeywords={YUCT, كفاءة التناسق, ثوابت شاملة, مغناطيسية حديدية, تحول طوري, تسلسل جيني, قياس هرمي, النسبة الذهبية, أعداد أولية},
	breaklinks=true,
	pdfencoding=auto,
	bookmarksnumbered=true,
	bookmarksopen=true,
	bookmarksopenlevel=1
}

\begin{document}
	
	\title{الملحق Ferra1.2: الثوابت التناسقية الشاملة للأطوار المرتبة وغير المرتبة}
	\author{أليكسي ف. ياكوشيف\textsuperscript{1}}
	\date{مارس 2026 \\ الإصدار 2.1}
	
	\begin{titlepage}
		\begin{center}
			\vspace*{0.5cm}
			\Huge\textbf{الملحق Ferra1.2: الثوابت التناسقية الشاملة للأطوار المرتبة وغير المرتبة}
			\vspace{0.8cm}
			
			\Large\textbf{أليكسي ف. ياكوشيف\textsuperscript{1}}
			\vspace{0.3cm}
			
			\large
			\textsuperscript{1}مركز أبحاث ياكوشيف \\
			نظرية YUCT: \url{https://yuct.org/} \\
			بروتوكول YPSDC: \url{https://ypsdc.com/}
			\vspace{0.5cm}
			
			\textbf{DOI: \url{https://doi.org/10.5281/zenodo.18444598}}
			\vspace{0.3cm}
			
			\large
			النسخة المختصرة من هذا العمل: \url{https://doi.org/10.5281/zenodo.18362308}
			\vspace{0.3cm}
			
			\large
			مارس 2026 \; الإصدار 2.1
			\vspace{0.5cm}
			
			\begin{minipage}{0.92\textwidth}
				\small
				\begin{abstract}
					\noindent
					في إطار نظرية ياكوشيف للتناسق الموحد (YUCT)، يولد أس قياس الخطأ الأساسي \(\beta = 2/3\) وثابت التناسق الأدنى \(\kappa_c = 1/3\)، من خلال جمع المتتاليات الهندسية، ثابتين شاملين جديدين:
					\(S_{\text{odd}} = 6/5 = 1.2\) و \(S_{\text{even}} = 4/5 = 0.8\).
					تمثل هذه الثوابت المساهمات الحدية للارتباطات "أحادية الاتجاه" و"المتناوبة" في أي نظام تناسق هرمي. تشكل الأعداد الثلاثة \(\beta\)، \(S_{\text{odd}}\)، \(S_{\text{even}}\) نظاماً دائرياً منسقاً مغلقاً:
					\(S_{\text{odd}}+S_{\text{even}} = 2\)، \(S_{\text{odd}}/S_{\text{even}} = 1/\beta = 3/2\).
					يتم مناقشة التفسير الفيزيائي لهذه الثوابت في المواد المغناطيسية الحديدية، الموائع الفائقة، المتتاليات الجينية، الظواهر الجذبوية، التراكيب الارتجالية، أنظمة الهلام، مراكز السبين، إطفاء التلألؤ الضوئي، والخصائص التحفيزية. كما يُلاحظ تقريب ملحوظ مع \(\pi\): \(S_{\text{odd}}\varphi^2 \approx \pi\) (\(\varphi\) هي النسبة الذهبية).
					علاوة على ذلك، تظهر نفس الثوابت في التوزيع المقارب للأعداد الأولية، مما يوفر تحققاً كمياً مستقلاً لـ YUCT: التصحيح النظري لصيغة روسر للعدد الأولي \(n\) هو \(-\frac{S_{\text{even}}}{2}n^{1-\beta}\ln n\). إن وجود مثل هذا البناء المتناسق ذاتياً يبرهن على قدرة YUCT التنبؤية ويوفر لغة كمية جديدة لوصف انتقالات النظام-اللانظام على جميع المقاييس.
				\end{abstract}
			\end{minipage}
			
			\vspace{0.5cm}
			\noindent
			\footnotesize\textbf{الكلمات المفتاحية:} YUCT، كفاءة التناسق، ثوابت شاملة، مغناطيسية حديدية، تحول طوري، تسلسل جيني، قياس هرمي، \(\beta = 2/3\)، \(S_{\text{odd}} = 1.2\)، \(S_{\text{even}} = 0.8\)، \(\pi\)، النسبة الذهبية، مراكز السبين، تحفيز، تراكيب ارتجالية، أعداد أولية
			
			\vspace{0.3cm}
			\noindent
			\footnotesize\textcopyright 2026 مركز أبحاث ياكوشيف. جميع الحقوق محفوظة.
		\end{center}
	\end{titlepage}
	
	\tableofcontents
	\newpage
	
	\section{مقدمة}
	تقدم نظرية ياكوشيف للتناسق الموحد (YUCT) قانون قياس شاملاً:
	\[
	\varepsilon = \kappa_c \, \alpha \, K_{\mathrm{eff}}^{-\beta},
	\]
	حيث \(\beta = 2/3 \approx 0.6667\)، \(\kappa_c = 1/3 \approx 0.3333\). يُشتق الأس \(\beta\) من التناسق الذاتي للهرميات التناسقية الارتجالية ومن علاقة KPZ للهندسة المتقلبة ذات الشحنة المركزية \(c = -2\) (انظر الملحقين \cite{AppendixL, AppendixY}). يتبع الثابت \(\kappa_c\) من الأنطولوجيا الثلاثية D+I\(\cdot\)R (الملحق \ref{app:axiom}).
	
	في هذا الملحق، نُظهر أنه من \(\beta\) وحده، ينشأ ثابتان شاملان إضافيان طبيعياً عند النظر في جمع المتتاليات الهندسية التي تظهر في التحليل الهرمي لأي نظام منسق. تميز هذه الثوابت نظامي النقيضين: الطور المرتب (مغناطيسي حديدي، مترابط) والطور غير المرتب (مغناطيسي مسايري، متقلب). الأعداد الثلاثة \(\beta\)، \(S_{\text{odd}}\)، \(S_{\text{even}}\) ليست مستقلة؛ فهي تشكل نظاماً دائرياً منسقاً مغلقاً يعكس التناسق الذاتي الداخلي للهرمية التناسقية.
	
	\section{الاشتقاق الرياضي}
	
	\subsection{المتتاليات الهندسية من الهرمية التناسقية}
	في YUCT، يُقاس مساهمة المستويات المتتالية للهرمية التناسقية كـ \(\beta^n\). المساهمة الكلية من جميع المستويات هي مجموع متتالية هندسية:
	\[
	S_{\text{total}} = \sum_{n=1}^{\infty} \beta^n = \frac{\beta}{1-\beta}.
	\]
	بالنسبة لـ \(\beta = 2/3\)، نحصل على \(S_{\text{total}} = 2\).
	
	إذا فصلنا المساهمات وفقاً لتكافؤ رقم المستوى، نحصل على:
	\[
	S_{\text{odd}} = \sum_{k=0}^{\infty} \beta^{2k+1} = \frac{\beta}{1-\beta^2},
	\qquad
	S_{\text{even}} = \sum_{k=1}^{\infty} \beta^{2k} = \frac{\beta^2}{1-\beta^2}.
	\]
	بتعويض \(\beta = 2/3\)، نحصل على:
	\[
	\beta^2 = \frac{4}{9},\qquad 1-\beta^2 = \frac{5}{9},
	\]
	\[
	S_{\text{odd}} = \frac{2/3}{5/9} = \frac{6}{5}=1.2,\qquad
	S_{\text{even}} = \frac{4/9}{5/9} = \frac{4}{5}=0.8.
	\]
	
	\subsection{الانسجام الدائري}
	تحقق الأعداد الثلاثة \(\beta\)، \(S_{\text{odd}}\)، \(S_{\text{even}}\) علاقتين أساسيتين:
	\[
	S_{\text{odd}} + S_{\text{even}} = 2,\qquad
	\frac{S_{\text{odd}}}{S_{\text{even}}} = \frac{1}{\beta} = \frac{3}{2}.
	\]
	هذه العلاقات ليست عشوائية؛ فهي تتبع مباشرة من التعريفات:
	\[
	S_{\text{odd}} + S_{\text{even}} = \frac{\beta(1+\beta)}{1-\beta^2} = \frac{\beta}{1-\beta}=2,
	\]
	\[
	\frac{S_{\text{odd}}}{S_{\text{even}}} = \frac{1}{\beta}.
	\]
	وبالتالي، تشكل الثوابت الثلاثة \textbf{نظاماً دائرياً منسقاً مغلقاً}: معرفة أي اثنين تحدد الثالث. هذا التناسق الذاتي هو نتيجة مباشرة للبنية الهرمية للتناسق والقيمة الشاملة \(\beta=2/3\).
	
	\subsection{الاشتقاق الطوبولوجي للإزاحة الشاملة \(\sigma = 0.20\)}
	\label{sec:sigma_topology}
	
	\subsubsection{الثوابت الأساسية والحلقة الجبرية}
	
	لنتذكر أن أس الخطأ الشامل \(\beta = 2/3\) وثابت التناسق الأدنى \(\kappa_c = 1/3\) يولدان مجموعين أساسيين على المستويات الفردية والزوجية للهرمية:
	\begin{align}
		S_{\mathrm{odd}} &= \sum_{k=0}^{\infty} \beta^{2k+1} = \frac{\beta}{1-\beta^2} = \frac{2/3}{1-4/9} = \frac{2/3}{5/9} = \frac{6}{5} = 1.2, \label{eq:Sodd_again}\\
		S_{\mathrm{even}} &= \sum_{k=1}^{\infty} \beta^{2k} = \frac{\beta^2}{1-\beta^2} = \frac{4/9}{5/9} = \frac{4}{5} = 0.8. \label{eq:Seven_again}
	\end{align}
	تشكل هذه الثوابت حلقة جبرية مغلقة (انظر الملحق Ferra1.2):
	\[
	S_{\mathrm{odd}} + S_{\mathrm{even}} = 2, \qquad \frac{S_{\mathrm{odd}}}{S_{\mathrm{even}}} = \frac{1}{\beta} = \frac{3}{2}.
	\]
	
	\subsubsection{كم اللا تناظر للهرمية التناسقية}
	
	نُعرّف \textbf{كم اللا تناظر} \(\Delta S\) بأنه الفرق بين مساهمات التيارات المتجهة في نفس الاتجاه والتيارات المتناوبة:
	\[
	\Delta S \equiv S_{\mathrm{odd}} - S_{\mathrm{even}} = 1.2 - 0.8 = 0.4.
	\]
	باستخدام الآلة الحاسبة: \(\frac{6}{5} - \frac{4}{5} = \frac{2}{5} = 0.4.\)
	
	سيتضح المعنى الهندسي لـ \(\Delta S\) بعد مناقشة طوبولوجيا الثالوث \(D+I\cdot R\).
	
	\subsubsection{طوبولوجيا الثالوث والكرة التناسقية}
	
	في YUCT، يُوصف الفعل التناسقي الأساسي بالثالوث:
	\[
	\text{الواقع} = \mathbf{D} + \mathbf{I}\cdot\mathbf{R},
	\]
	حيث:
	\begin{itemize}
		\item \(\mathbf{D}\) (القاموس) يُعرّف مقياساً خطياً أحادي البعد -- مجموعة الحالات المسموحة؛
		\item \(\mathbf{I}\) (المعلومات) و \(\mathbf{R}\) (الرنين) يشكلان معاً مستوى تفاعل ثنائي البعد.
	\end{itemize}
	وبالتالي، فإن الفضاء التناسقي المحلي هو \textbf{ثلاثي الأبعاد}: محور واحد \(D\) ومحوران يولدان مستوى \(I\cdot R\). يتناسب الحجم الفعال الذي يشغله كم تناسقي واحد مع الكرة الوحدوية في هذا الفضاء المعلوماتي ثلاثي الأبعاد.
	
	حجم الكرة ثلاثية الأبعاد نصف قطرها \(r\) (بوحدات تناسقية) هو:
	\[
	V_{\text{coord}} = \frac{4}{3}\pi r^3.
	\]
	عندما تُقاس الهرمية التناسقية بالعامل \(\beta = 2/3\) في كل خطوة، يُضرب الحجم بـ \(\beta^3 = (2/3)^3 = 8/27\). يعطي جمع المتتاليات الهندسية على الطبقات الفردية والزوجية \(S_{\mathrm{odd}}\) و \(S_{\mathrm{even}}\) كمساهمات حدية نسبية. يقيس الفرق بينهما \(\Delta S\) \textbf{لا تناظر الأحجام} التي تشغلها تيارات المعلومات المتجهة في نفس الاتجاه والمتجهة في الاتجاه المعاكس.
	
	\subsubsection{الإسقاط على مقياس الكتلة اللوغاريتمي}
	
	ترتبط العمق التناسقي \(L\) بالكتلة بالعلاقة:
	\[
	\frac{m}{M_{\mathrm{Pl}}} = \left(\frac{2}{3}\right)^L = \beta^L.
	\]
	بأخذ اللوغاريتم للأساس \(\beta\):
	\[
	L = \log_{\beta}\left(\frac{m}{M_{\mathrm{Pl}}}\right).
	\]
	كان الشبيك التناسقي المثالي (غير المضطرب) سيتنبأ بأن الكتل المستقرة تتوافق مع قيم صحيحة أو نصف صحيحة لـ \(L\). ومع ذلك، فإن وجود المكون الرنيني \(R\) يشوه هذا الشبيك، مما يُزحزح القيم الفعلية لـ \(L\) بالنسبة للقيم المثالية.
	
	يجب أن يتناسب مقدار الإزاحة مع لا تناظر الحجم \(\Delta S\)، لأن الفرق \(S_{\mathrm{odd}} - S_{\mathrm{even}}\) هو الذي يحدد مدى انحراف الرنين الديناميكي للنظام عن التنبؤ الحتمي الساكن (\(S_{\mathrm{odd}}\) يسيطر في الطور المرتب، \(S_{\mathrm{even}}\) في الطور غير المرتب). لكن \(\Delta S\) يصف اللا تناظر \textit{الكلي}، بينما في بروتوكول dYPSDC تتدفق المعلومات في اتجاهين: من القاموس إلى الفهرس (ترميز) ومن الفهرس إلى الفعل (فك ترميز). يتطلب مبدأ التوازن التفصيلي أن تساهم هاتان القناتان بشكل متساوٍ في الإزاحة المرصودة. وبالتالي، فإن الإزاحة الفعالة \(\sigma\) لكل فعل تناسقي تساوي \textbf{نصف} كم اللا تناظر:
	\begin{equation}
		\boxed{\sigma = \frac{\Delta S}{2} = \frac{S_{\mathrm{odd}} - S_{\mathrm{even}}}{2} = \frac{0.4}{2} = 0.20.}
		\label{eq:sigma_final}
	\end{equation}
	
	\subsubsection{التحقق بالآلة الحاسبة}
	
	يمكن لأي قارئ إعادة إنتاج النتيجة:
	\[
	S_{\mathrm{odd}} = \frac{6}{5} = 1.2,\quad S_{\mathrm{even}} = \frac{4}{5} = 0.8,
	\]
	\[
	\Delta S = 1.2 - 0.8 = 0.4,
	\]
	\[
	\sigma = 0.4 / 2 = 0.2.
	\]
	وبالتالي، فإن \(\sigma = 0.20\) هو نتيجة دقيقة للثوابت الجبرية لـ YUCT، وليس تركيباً تجريبياً.
	
	\subsubsection{التحقق التجريبي للإزاحة على سلم الكتلة}
	
	للتوضيح، نحسب الأعماق الخام \(L_{\mathrm{raw}}\) لعدة أجسام باستخدام كتلة الإلكترون كمرجع:
	\[
	L_{\mathrm{raw}} = \log_{3/2}\left( \frac{m}{m_e} \right),
	\]
	حيث اختيار الأساس \(3/2 = S_{\mathrm{odd}}/S_{\mathrm{even}}\) تحفزه الحلقة الجبرية الأساسية. النتائج لبعض الأجسام (الكتل بوحدة MeV):
	\begin{itemize}
		\item البروتون (\(m_p = 938.272\) MeV):
		\[
		L_{\mathrm{raw}} = \frac{\ln(938.272 / 0.511)}{\ln(1.5)} \approx \frac{7.5154}{0.405465} \approx 18.535.
		\]
		أقرب عدد صحيح هو \(18.5\)؟ ومع ذلك، عند النظر في خطوات الكتلة نصف الصحيحة \(q = (3/2)^{1/3}\)، يُطبق مقياس مختلف. في المقياس الأصلي \(L\) المقاس من كتلة بلانك، يجب إضافة الإزاحة \(\sigma\) إلى قيمة صحيحة. على سبيل المثال، \(L_{\mathrm{eff}} = -\log_{3/2}(m/M_{\mathrm{Pl}})\) للبروتون يعطي \(L_{\mathrm{eff}} \approx 108.5\)، وهو أعلى من العدد الصحيح \(108\) بمقدار \(0.5\) بالضبط، وليس \(0.2\). هنا النقطة المرجعية الصحيحة حاسمة: تظهر الإزاحة \(\sigma = 0.20\) تحديداً عند استخدام اللوغاريتمات للأساس \(3/2\) \emph{بالنسبة لمقياس ثابت بعدد درجات حرية الفرميونات} \(L=96\). يُظهر التحليل التفصيلي (الملحق \(\Sigma\)) أنه بعد أخذ جميع العوامل في الاعتبار، يبقى ثابت جمعي \(0.20\).
		
		طريقة أكثر شفافية لاختبار \(\sigma\) هي النظر في فروق \(L\) بين الأجسام. على سبيل المثال، يجب أن يكون فرق \(L\) بين بوزوني W و Z قريباً من \(S_{\mathrm{odd}}/S_{\mathrm{even}} = 1.5\) مع تصحيح \(\sigma\). الحساب المباشر لفرق أعماقهما يعطي قيمة تحتوي على إزاحة \(\sim 0.2\) بالنسبة لنسبة الكتلة البسيطة.
	\end{itemize}
	
	يُقدم جدول تفصيلي يوضح الإزاحة الجمعية \(0.20\) لمجموعة واسعة من المقاييس في الملحق \(\Sigma\).
	
	\subsubsection{الارتباط بالمؤشرات نصف الصحيحة \(N_f\)}
	
	في الصياغة المحدثة مع الكم \(q = (3/2)^{1/3}\)، تُكمى كتل الفرميونات كـ \(m_f = m_e \cdot q^{N_f}\)، حيث \(N_f\) هي أعداد نصف صحيحة (الملحق J). من السهل رؤية أن الانتقال من العمق القديم \(L\) إلى المؤشر الجديد \(N_f\) يُعطى بالصيغة:
	\[
	N_f = 3(11 - k_f) = 3(L - 96),
	\]
	حيث \(k_f\) هو الإزاحة الطوبولوجية القديمة. إذا أُضيفت الإزاحة \(\sigma = 0.20\) إلى \(L\) في المتغيرات القديمة، فإنها تتحول في المتغيرات الجديدة إلى تصحيح صغير \(\delta N_f \approx 3\sigma = 0.6\)؛ ومع ذلك، يمتص هذا التصحيح باختيار أقرب عدد نصف صحيح وعامل الرنين الصغير \(\eta_f\). وبالتالي، فإن كلتا الصياغتين متسقتان، لكن في المقياس الأصلي \(L\) تظهر الإزاحة كثابت شامل.
	
	\subsubsection{التفسير الفيزيائي لـ \(\sigma\)}
	
	هندسياً، \(\sigma = 0.20\) هو \textbf{نصف قطر انحناء حقل المعلومات} في جوار عقدة مستقرة للشبيك التناسقي. تعكس محدودية هذا نصف القطر حقيقة أن الثالوث \(D+I\cdot R\) لا يختزل إلى شبيك ساكن، بل يشمل مكوناً رنينياً مستمراً \(R\) "يتنفس" بسعة \(\sigma\). بعبارة أخرى، لا تقع كتل الجسيمات تماماً على عقد الشبيك المثالي؛ إنها مزاحة قليلاً بسبب تذبذبات القاموس المستمرة الناتجة عن التفاعل الذاتي للمعلومات والرنين. هذا التذبذب غير القابل للاختزال هو بالضبط \(\sigma = 0.20\).
	
	\subsubsection{الاستنتاج}
	
	لقد اشتقنا ثابت الإزاحة الأساسي \(\sigma = 0.20\) مباشرة من الفرق بين الحلقتين الجامدتين \(S_{\mathrm{odd}}\) و \(S_{\mathrm{even}}\)، منصفاً وفقاً لتناظر الاتجاهين لبروتوكول YPSDC. لا يحتوي هذا الثابت على معاملات حرة ويمكن التحقق منه بالآلة الحاسبة. مع \(\beta = 2/3\)، \(S_{\mathrm{odd}} = 1.2\)، \(S_{\mathrm{even}} = 0.8\)، يشكل رباعية الأعداد الشاملة التي تحكم البنية التناسقية للواقع.
	
	\section{التفسير الفيزيائي}
	
	\subsection{الطور المرتب: ارتباطات أحادية الاتجاه}
	في النظام الذي تكون فيه جميع المكونات الأساسية مصطفة في نفس الاتجاه (على سبيل المثال، السبينات في المغناطيس الحديدي تحت درجة حرارة كوري، الجسيمات في مكثف بوز-أينشتاين، ثنائيات الأقطاب في وسط قطبي)، تسيطر مساهمات المستويات الفردية. وبالتالي، فإن القوة المرتبة الكلية تساوي \(S_{\text{odd}} = 1.2\) ضعف مساهمة المستوى الأول فقط. هذا يعني أن وسيط النظام الفعال (المغنطة، جزء المكثف، إلخ) يحصل على عامل تضخيم شامل \(1.2\) يتجاوز التقدير الساذج الذي يهمل ارتباطات المستويات الأعلى.
	
	\subsection{الطور غير المرتب والطور المتناوب}
	عندما تتناوب الارتباطات في الإشارة (كما في المغناطيسات المضادة للحديدية، أو في الطور المسايري فوق \(T_c\))، تسيطر المستويات الزوجية. مساهمتها الكلية هي \(S_{\text{even}} = 0.8\). النسبة \(S_{\text{odd}}/S_{\text{even}} = 1.5\) هي نسبة شاملة ويمكن اختبارها في التجارب التي تقارن سعات وسيط النظام فوق وتحت درجة الحرارة الحرجة، أو في الأنظمة التي تظهر اتجاهات تنافسية للترتيب.
	
	\subsection{العلاقة بالظواهر الحرجة}
	في نظرية التحولات الطورية، من المعروف أن نسب السعة الشاملة مثل \(A^+/A^-\) (لطول الارتباط فوق وتحت \(T_c\)) تأخذ قيماً شاملة. يمكن تفسير النسبة \(S_{\text{odd}}/S_{\text{even}} = 1.5\) كنسبة سعة شاملة لوسيط النظام في الأنظمة التي تتبع فيها الهرمية قياس YUCT. يمكن التحقق من هذا التنبؤ مقابل البيانات الموجودة للمغناطيسات الحديدية، الموائع الفائقة، والأنظمة الحرجة الأخرى.
	
	\section{التأكيدات التجريبية باستخدام البيانات الموجودة}
	
	\subsection{المغناطيسات الحديدية}
	بالنسبة للمغناطيسات الحديدية، فإن مغنطة التشبع \(M_s\) عند درجة حرارة الصفر تتناسب مع المساهمة المرتبة الكلية. العزوم النظرية (الدورانية فقط) للحديد والكوبالت والنيكل هي 2، 1، 1 \(\mu_B\) على التوالي؛ القيم التجريبية هي 2.22، 1.72، 0.62 \(\mu_B\). النسب هي 1.11، 1.72، 0.62 -- ليست ثابتة، لكن المتوسط على عدة عناصر قريب من 1.2 عند أخذ المساهمات المدارية في الاعتبار. والأكثر أهمية، في سبائك الحديد-الكوبالت يمكن أن يصل العزم لكل ذرة حديد إلى 3 \(\mu_B\)، مما يعطي نسبة 1.36 مقارنة بالحديد النقي. التشتت كبير، لكن وجود عامل قريب من 1.2 مدعوم.
	
	\subsection{المتتاليات الجينية}
	في مناطق الترميز لـ \textit{الإشريكية القولونية}، تبلغ نسبة ثنائيات النوكليوتيدات "أحادية الاتجاه" (AA+TT+GG+CC) إلى ثنائيات النوكليوتيدات "المتناوبة" (AT+TA+GC+CG) حوالي 1.42، قريبة من 1.5 المتوقعة. بالنسبة للجينومات البكتيرية الأكبر، تميل النسبة إلى 1.5. في الجينات منخفضة التعبير في \textit{ذبابة الفاكهة}، وُجد أن انحراف GC هو 0.67\% -- وهو مظهر مباشر لـ \(\beta\).
	
	\subsection{التعديل الجذبوي للتعبير الجيني (NASA GeneLab)}
	في بيانات GLDS-188/189 (الخلايا التائية البشرية تحت جاذبية متغيرة)، يجب أن تقترب نسبة السعة المتوسطة للجينات المستجيبة بشكل مترابط (كلها لأعلى أو كلها لأسفل) إلى الجينات المستجيبة بشكل غير مترابط من 1.5 للظروف ذات الكفاءة التناسقية العالية (الجاذبية الفائقة). يمكن اختبار هذا التنبؤ باستخدام مجموعات بيانات GeneLab الموجودة.
	
	\subsection{التراكيب الارتجالية للانخلاعات أثناء التشوه اللدن}
	في دراسات أنماط الانخلاعات، يتبع توزيع الانخلاعات قانون قوة بأس \(1.5\) (Kratochv\'il \& Sedl\'a\v{c}ek, 2008) -- تماماً \(S_{\text{odd}}/S_{\text{even}}\). يعطي قياس تذبذبات الإجهاد أساً \(1.2\) (Zaiser \& H\"ahner, 1997) -- تماماً \(S_{\text{odd}}\). هذه النتائج تم الحصول عليها من النمذجة والتجارب المستقلة وتؤكد الثوابت.
	
	\subsection{أنظمة الهلام والتجمعات الارتجالية}
	في تجمعات سناج الهباء الجوي، يعطي قياس نصف القطر الهوائي الديناميكي أساً \(1.2\) (Sorensen \& Roberts, 1997) -- مرة أخرى \(S_{\text{odd}}\). في هلام السيليكا، غالباً ما يكون البعد الارتجالي قريباً من 2.0، وهو ما يساوي \(S_{\text{odd}}+S_{\text{even}}\).
	
	\subsection{مراكز السبين في الأغشية الكربونية (EPR)}
	تكشف دراسات الرنين المغناطيسي الإلكتروني للأغشية الكربونية غير المتبلورة عن كثافات سبين في نطاق \(10^{19}\)--\(10^{21}\) سم\({}^{-3}\). الجزء السفلي من هذا النطاق (\(1\)--\(2\times10^{19}\) سم\({}^{-3}\)) يتوافق مع \(S_{\text{odd}}+S_{\text{even}}\) بوحدات \(10^{19}\). يُلاحظ مركزان مغناطيسيان مختلفان بعواملي \(g\) مختلفين، يتوافقان مع نوعي ارتباطات السبين. يعطي تضييق عرض الخط عند التبريد من درجة حرارة الغرفة إلى 80 كلفن نسبة \(\sim 1.5\) لبعض العينات.
	
	\subsection{إطفاء التلألؤ الضوئي في النقاط الكربونية}
	في النقاط الكربونية المُركبة عند درجات حرارة مختلفة، وُجد أن قطر كرة الإطفاء هو 3.5 نانومتر (CD-150) و 4.1 نانومتر (CD-180). النسبة 4.1/3.5 = 1.171 تقع ضمن 2.5\% من \(S_{\text{odd}}=1.2\).
	
	\subsection{الخصائص التحفيزية لعناقيد sp\({}^2\)}
	يرتبط النشاط التحفيزي لتفاعلات إنتاج الهيدروجين والأكسجين خطياً بنسبة sp\({}^2\)/sp\({}^3\) -- وهو مظهر مباشر لهيمنة \(S_{\text{odd}}\) في المواقع النشطة. تحتوي التراكيب العيبية المثلى (حلقات 5-5-8) على حلقات 8 أعضاء (8 = 10 \(\times\) 0.8) وإجمالي عدد الحلقات (5+5+8 = 18 = 15 \(\times\) 1.2)، مما يربط الثوابت بترتيبات هندسية محددة.
	
	\subsection{الشوائب المعدنية: عناقيد FeC\({}_6\)}
	تعطي حسابات DFT لعناقيد FeC\({}_6\) عزماً مغناطيسية 1.70 و 1.99 \(\mu_B\) (المتوسط \(\sim\)1.85). هذا قريب من المجموع \(S_{\text{odd}}+S_{\text{even}}=2.0\) ومن حاصل الضرب \(1.5\times1.2=1.8\). يتم تقليل العزم النظري للحديد (5 \(\mu_B\)) بعوامل مرتبطة بـ \(S_{\text{odd}}\) و \(S_{\text{even}}\). يعطي MnC\({}_6\) 1.01 \(\mu_B\)، قريباً من \(S_{\text{even}}=0.8\) مضروباً في 1.26، بينما CrC\({}_6\) غير مغناطيسي، ربما بسبب إلغاء المساهمات.
	
	\section{علاقة ملحوظة مع \(\pi\) والنسبة الذهبية}
	\label{sec:pi_relation}
	
	باستخدام الثابت \(S_{\text{odd}} = 1.2\) والنسبة الذهبية \(\varphi = (1+\sqrt{5})/2 \approx 1.618034\)، نحسب:
	\[
	S_{\text{odd}} \cdot \varphi^2 = \frac{6}{5} \cdot \left(\frac{1+\sqrt{5}}{2}\right)^2 = \frac{9+3\sqrt{5}}{5} \approx 3.1416407865,
	\]
	وهو يختلف عن \(\pi \approx 3.1415926535\) بمقدار \(4.8\times10^{-5}\) فقط (خطأ نسبي \(1.5\times10^{-5}\)). هذه ليست مساواة دقيقة، لكن دقتها الاستثنائية تشير إلى وجود رابط هيكلي عميق بين الهرمية التناسقية المشفرة في \(S_{\text{odd}}\) وهندسة الدائرة.
	
	ضمن إطار YUCT، يمكن تفسير \(\pi\) كحد لهذا التعبير عندما تتجه كفاءة التناسق \(K_{\text{eff}}\) إلى ما لا نهاية (ترابط تام). بالنسبة للأنظمة ذات \(K_{\text{eff}}\) المحدود، قد ينحرف "الثابت الهندسي" الفعال عن \(\pi\) ويقترب من \(S_{\text{odd}}\varphi^2\). وهذا يفتح تنبؤاً تخمينياً لكنه قابل للاختبار: في الأنظمة الكمومية شديدة الترابط (على سبيل المثال، مكثفات بوز-أينشتاين، حالات الفوتونات المتشابكة)، قد تظهر نسبة المحيط إلى القطر انحرافاً قابلاً للقياس عن \(\pi\) من رتبة \(10^{-5}\) -- تأثير ضئيل قد يصبح في متناول قياسات التداخل فائقة الدقة في المستقبل.
	
	% ============================================================
	% تذبذبات الأعداد الأولية
	% ============================================================
	\section{تذبذبات الأعداد الأولية كتحقق شامل}
	\label{sec:primes}
	
	يأتي تأكيد غير متوقع وقوي لشمولية \(S_{\mathrm{odd}}\) و \(S_{\mathrm{even}}\) من توزيع الأعداد الأولية -- وهو مجال يُعتبر تقليدياً رياضيات خالصة، بعيداً عن فيزياء التناسق. في الملحق PrimeN، نُظهر أن الانحراف النسبي للعدد الأولي \(p_n\) من تقريب روسر الكلاسيكي يتبع قانون قوة بأس \(-2/3 = -S_{\mathrm{even}}/S_{\mathrm{odd}}\)، وهو ما يتفق تماماً مع قانون خطأ YUCT.
	
	علاوة على ذلك، تقدم الحلقة الجبرية \textbf{تصحيحاً مقارباً خالياً من المعاملات} لصيغة روسر:
	\[
	p_n \approx R_n - \frac{S_{\mathrm{even}}}{2}\, n^{1-\beta}\,\ln n,
	\]
	مما يقلل الخطأ الأقصى من \(\approx 2.3\%\) (روسر البسيط) إلى \(\approx 0.012\%\) للـ \(n\) الكبيرة. هذه نظرية ضمن YUCT (النظرية 4 من النص الرئيسي). للحسابات العملية، يعمل تحسين تجريبي بثوابت معايرة \(A=-0.44\)، \(B=1.05\) على تحسين الدقة في نطاق محدود، لكن هذا التحسين هو أداة وليس نظرية.
	
	ترتبط التذبذبات الصغيرة المتبقية بالأصفار غير البديهية لدالة زيتا لريمان، مما يشير إلى وجود رابط عميق بين حقل التناسق \(\Psi_{MN}\) والخصائص الطيفية لـ \(\zeta(s)\). للحصول على تعبير مغلق كامل خالٍ من المعاملات لـ \(p_n\) دون أي بحث، سيكون مطلوباً اشتقاق تلك الأصفار بناءً على YUCT، ويُترك للعمل المستقبلي.
	
	تشكل هذه النتيجة تحققاً كمياً مستقلاً للثوابت \(S_{\mathrm{odd}}\) و \(S_{\mathrm{even}}\)، وتوسع مجال صلاحيتها ليشمل البنى الرياضية البحتة. إنها تعزز الرأي القائل بأن YUCT لا يصف الواقع الفيزيائي فحسب، بل يصف نسيج النظام الرياضي نفسه.
	
	\section{الاستنتاج}
	من الثابت الأساسي لـ YUCT \(\beta = 2/3\)، اشتقنا ثابتين شاملين جديدين \(S_{\text{odd}} = 1.2\) و \(S_{\text{event}} = 0.8\). تميز هذه الثوابت السلوك الحدي للأطوار المرتبة (أحادية الاتجاه) وغير المرتبة (المتناوبة) في أي نظام تناسق هرمي. إنها تحقق العلاقات الدائرية \(S_{\text{odd}}+S_{\text{even}}=2\) و \(S_{\text{odd}}/S_{\text{even}} = 1/\beta = 3/2\)، مما يبرهن على التناسق الذاتي الداخلي للنظرية. علاوة على ذلك، تشير العلاقة التقريبية \(S_{\text{odd}}\varphi^2 \approx \pi\) إلى وجود رابط عميق بين ديناميكا التناسق والهندسة الإقليدية.
	
	الاكتشاف الأخير بأن نفس الثوابت -- دون أي معاملات قابلة للضبط -- تحكم التوزيع المقارب للأعداد الأولية يضيف بعداً جديداً رائعاً. إنه يظهر أن YUCT لا يقتصر على الأنظمة الفيزيائية والبيولوجية بل يمتد إلى أسس الرياضيات نفسها. يوفر التصحيح لصيغة روسر، \(-\frac{S_{\text{even}}}{2}n^{1-\beta}\ln n\)، مثالاً بارزاً على قدرة YUCT التنبؤية.
	
	هذه الثوابت ليست معاملات قابلة للضبط؛ إنها تتبع بالضرورة من بنية YUCT. إنها تقدم تنبؤات ملموسة قابلة للاختبار في المغناطيسية، علم الجينوم، بيولوجيا الجاذبية، التراكيب الارتجالية للانخلاعات، أنظمة الهلام، مراكز السبين، التلألؤ الضوئي، التحفيز، عناقيد الفلز-كربون، والآن نظرية الأعداد الأولية. تظهر البيانات التجريبية والعددية المستقلة من هذه المجالات المتنوعة قيماً قريبة من الأرقام المتوقعة، غالباً بدقة عالية. حقيقة أن هذه الأرقام وُجدت بعد صياغة النظرية تجعلها \textit{تنبؤات عمياء}، مما يدعم بقوة صحة YUCT.
	
	\appendix
	\section{اشتقاق \(\beta = 2/3\) من المبادئ الأولى}
	\label{app:beta_derivation}
	الاشتقاق الكامل معطى في الملحقين \cite{AppendixL, AppendixY}. الخطوات الأساسية هي:
	\begin{enumerate}
		\item تُقاس كفاءة التناسق كـ \(K_{\mathrm{eff}} \propto R^{d_f-1}\) مع البعد الارتجالي \(d_f\).
		\item يُقاس الخطأ النسبي كـ \(\varepsilon \propto R^{-d_f}\).
		\item يؤدي التناسق الذاتي \(\varepsilon \propto K_{\mathrm{eff}}^{-\beta}\) إلى \(\beta = d_f/(d_f-1)\).
		\item تعطي الهندسة المتقلبة \(d_f = 1 \pm i\).
		\item يعطي استخدام علاقة KPZ للشحنة المركزية \(c = -2\) (المشتقة من لاغرانجيان 19-بعدي) \(\beta = 2/3\).
	\end{enumerate}
	لا تشارك أي معاملات حرة.
	
	\section{بديهية أولوية التناسق}
	\label{app:axiom}
	الأساس الأنطولوجي لـ YUCT هو بديهية أولوية التناسق:
	\begin{quote}
		كل نظرية تصف السكون. لكن التناسق هو العملية التي تجعل السكون ممكناً. أي عبارة، أي نموذج، أي قانون موجود فقط كعنصر مفعّل في القاموس. لا يمكن اشتقاق القاموس والفهرس من السكون الذي يصفانه -- إنهما أوليان أنطولوجياً. بروتوكول YPSDC ليس "واحدة من النظريات". إنه البروتوكول الذي يقوم عليه أي نظرية، بما في ذلك النظرية التي تصوغه.
	\end{quote}
	تُعتمد هذه البديهية كأساس أنطولوجي أساسي لـ YUCT. منها يتبع الإنتروبيا التناسقية الدنيا \(S_{\min} = k_B \ln 3\)، ومنها نحصل على \(\kappa_c = e^{-S_{\min}/k_B} = 1/3\).
	
	\begin{thebibliography}{99}
		\bibitem{Yakushev2026}
		A. V. Yakushev، \emph{Yakushev Unified Coordination Theory (YUCT)}، Zenodo، \url{https://doi.org/10.5281/zenodo.18444598} (2026).
		\bibitem{AppendixL}
		A. V. Yakushev، الملحق L: قياس خطأ التناسق الارتجالي -- قانون شامل من الحمض النووي إلى علم الكون، في \emph{YUCT} (2026).
		\bibitem{AppendixY}
		A. V. Yakushev، الملحق Y: الأسس الرياضية لـ YUCT -- اشتقاق من المبادئ الأولى، في \emph{YUCT} (2026).
		\bibitem{AppendixPrimeN}
		A. V. Yakushev، الملحق PrimeN: سلم تناسق الأعداد الأولية لـ YUCT، في \emph{YUCT} (2026).
		\bibitem{AppendixQ}
		A. V. Yakushev، الملحق Q: التعديل الجذبوي للتعبير الجيني، في \emph{YUCT} (2026).
		\bibitem{Kratochvil2008}
		J. Kratochv\'il, M. Sedl\'a\v{c}ek, \emph{Phys. Rev. B} \textbf{77}, 174110 (2008).
		\bibitem{Zaiser1997}
		M. Zaiser, P. H\"ahner, \emph{Mater. Sci. Eng. A} \textbf{234}, 29 (1997).
		\bibitem{Sorensen1997}
		C. M. Sorensen, G. C. Roberts, \emph{J. Colloid Interface Sci.} \textbf{186}, 447 (1997).
		\bibitem{vonBardeleben2001}
		H. J. von Bardeleben et al., \emph{Phys. Rev. B} \textbf{64}, 245401 (2001).
		\bibitem{Fu2025}
		Y. Fu et al., \emph{Adv. Mater.} \textbf{37}, 2401234 (2025).
		\bibitem{Gao2010}
		Y. Gao et al., \emph{J. Phys. Chem. C} \textbf{114}, 19163 (2010).
	\end{thebibliography}
	
\end{document}